\documentclass[11pt]{article}
\usepackage[margin=1in]{geometry}
\usepackage{amsmath}
\usepackage{booktabs}
\usepackage{hyperref}
\usepackage{graphicx}
\usepackage{soul}
\usepackage{xcolor}
\sethlcolor{yellow} % default is yellow


\title{A FOSS automated calibration software system for low-cost capacitive soil moisture sensors using ESP32 microcontrollers}
\author{Alex Southgate\\
\texttt{alex.southgate@equalitylabs.cymru}
}

\date{\today}

\begin{document}
\maketitle

\begin{abstract}

ESP32s are low-cost, low-energy microcontrollers with integrated bluetooth and WiFi, and have recently lowered the barrier for entry into prototyping smart devices. ESP32s equipped with capacitive moisture sensors have been suggested for low-cost smart plant watering systems. However, measuring moisture in soil is complex, potentially destructive, and requires careful calibration in order to characterise the response curve mapping soil water content to sensor measurements, which many users are unlikely to undertake. Even accurately measuring soil moisture content requires effort, with variability across soils and environmental conditions, and consumer-grade probes tending to offer coarse measurements only. In order for these low-cost smart sensor systems to be practically useful for this category of user, calibration must be easy, and uncertainty measures available. The objective of this work was to develop a near-automous, free and open source software (FOSS) system for monitoring, with continuously updating calibration and measures of uncertainty. 
\end{abstract}

\section{Introduction}

Recent work has proposed the use of ESP32-class microcontrollers and low-cost moisture sensors for automated watering systems and measurement of soil water content (SWC) \cite{biswas2018-solar-pump, purnomo2021-lora, sudarmaji2024-smart-soil, afzal2025-arima}, motivated by mitigating water shortages \cite{afzal2025-arima}, with increased interest in application of IoT sensors to agriculture \cite{ruiz2025-iot-mapping}. As global warming continuous, both flooding and water shortages are increasing world-wide \cite{ficklin2022hydrological}, including in the United Kingdom \cite{lane2021climate}. In previous work, ESP32/capacitive sensor calibration has been performed by gravimetry \cite{sudarmaji2024-smart-soil} and WiFi connections. However, it is expected that calibrated response curves depend on a range of factors, including soil type (i.e. clay content), and temperature \cite{merlin2007-calibration}. Furthermore, end users may find it difficult to undertake calibration for each specific application type. 

Typically, gravimetric sampling can be used to characterise SWC, but is destructive and time consuming; this process often involves drying soil by heating and therefore cannot be applied to existing plant systems. Other methods include time domain reflectometry and capacitance and impedence sensors \hl{cite}. Capacitive moisture sensors, as with all electromagnetic methods (EM) measure SWC indirectly through the dielectric permittivity \cite{kelleners2005-frequency-dependence}, since the dielectric permitivity of water exceeds that of other soil components \hl{cite}, though it is dependent on measurement electrical frequency \cite{kelleners2005-frequency-dependence,sheng2025-frequency} and temperature \cite{jones2005-em-methodology}. Field capacity (FC), the water content at which internal drainage stops, is often utilized \cite{jury2004soil}. 

Ideal ranges for SWC varies across plant species and growth stages. For example, tomatoes (\textit{Lycopersicon esculentum}) may grow best at 70\% field capacity (FC) \cite{liu2019tomato}, while the cactus \textit{Opuntia ficus-indica} was shown to exhibit stress under 30\% FC \cite{lahbouki2023effects}.\\

Existing work has been undertaken on moisture sensor calibration and characterisation of variability. Bogena et al (2017) \cite{bogena2017-calibration} developed a procedure for calibration of low-cost SMT1000 moisture sensors and undertook systematic calibration of hundreds of sensors in a range of media, observing sensor-to-sensor variability and both density and temperature dependence. In contrast, Kizito et al. (2008) \cite{kizito2008-frequency} found that for research-grade ECH2O probes (Decagon Devices Inc., Pullman, WA) and a measurement frequency of 70 MHz, a single calibration curve was sufficient for relating electrical conductivity to SWC across different soil salinities, though they did find variability with the degree of probe immersion and dependency of measurements on temperature, and the sensors used are more an order of magnitude more costly than the ``capacitive soil moisture sensor v.2.0'' used in other studies (e.g. \cite{sudarmaji2024-smart-soil}). 

The objective of this work was to introduce an automated system and statistical method for automated calibration of sensors with minimal user input. This aim was primarily motivated by the challenge end users may find with gravimetric calibration, as well as lack of generalisability of a given calibration curve. We aim to present a method providing automated measurements, calibration, and dashboards for visualisation and alerts.

\section{Methods}

\subsection{Software and hardware architecture}

\begin{figure}[t]
  \centering
  \includegraphics[width=1.0\linewidth]{images/hcult.pdf}
  \caption{Simplified software architecture. ESP32 devices repeat a cycle of deep-sleep, sample, advertise-only BLE. A monitor node continuously scans for advertising sensors, collects payloads, and writes them to the database. Occasionally, an inference daemon runs event detection and curve calibration. Finally, an API node allows users to run operations and retrieve data for dashboards. }
  \label{fig:myfigure}
\end{figure}

In our test system, ESP32s were used with low-cost capacitive moisture sensors. \hl{TODO}: description of boards, sensor sampling frequency, polling, etc.

\subsection{Segmentation}

A hierarchical approach was taken to segmentation. For the target system, three types of processes are defined at different time-scales: shocks caused by watering from above; smooth changes induced by watering from below; loss of moisture due to drying or uptake by a plant. All of these processes cause a response in sensor readings over a period of time. Given that the timescales of watering dynamics are much smaller than those of drying and uptake, we first defined `high-velocity' and `low-velocity' states, with the latter approximately corresponding to equilibrium when drying and uptake is slow. Firstly, irregularly spaced sensor data was smoothed with an exponential moving average \cite{genccay2001introduction} (or exponentially weighted moving average (EWMA) \cite{zumbach2001operators}), a single-pole low-pass filter \cite{smith1997scientist} with time-weighted coefficients:

\[
\tilde{X}_t = \exp^{-\Delta t/\tau} \tilde{X}_{i-1} + (1-\exp^{-\Delta t/\tau}) X_i
\]

After smoothing, the data was resampled and derivatives were approximated by finite differences (second order) \cite{griffiths2006numerical}. Two thresholds were defined for identification of high-velocity state regions: a trigger threshold $v_{trigger}$ and a release threshold $v_{release}$ (see relay with hysteresis \cite{khalil2002nonlinear}). Once the trigger is passed at a high instantaneous velocity, it must decrease beyond the release to return to the `low-velocity` state. \\

It is possible in operation that several watering events will happen in rapid succession before system equilibriation. So, given an initial segmented sequence of data into `high-velocity' (non-equilibrium) and `low-velocity' (close to equilibrium) regions, within-block characterisation was undertaken. In the high-velocity state, a second relay with hysteresis was applied to the second derivative.\\

\subsection{Watering event detection}

Raw sensor data is continuously streamed into the database system. From there, event detection is performed on recent time intervals. Event detection is performed in several steps: initial triggers based on outlier time series deltas; stable event detection by testing for stability of change in a window; subsetting to events with a sign consistent with watering; consensus amongst sensors for a single microcontroller; manual exclusion of known non-watering events such as sensor removal or repositioning; evaluation. Firstly, denoting measurement at $t_i$ of sensor $s$ as $X_{i,s}$, which are assumed to be continuously generated by the software system previously described, the distribution of $\Delta X_{i,s}$ in manually-defined event-free time intervals can be gathered, and with a threshold $p$, extreme values are taken as candidates if $|\Delta X_{i,s}|$ is greater than a threshold. In this case, using the median absolute deviation and compute modified Z-scores \cite{iglewicz1993volume}:

\[
M_{i,s} = \frac{0.6745 \cdot (|\Delta X_{i,s}| - median_{w}{|\Delta X_{\cdot,s}|})}{MAD_w}
\]

Given an integer lag $L$ to compute differences, a median window $w$, and a significance threshold $p$, an event was defined at a time point $t_i$ if $p_j > p$ for all $j$ between $i$ and $i+L$. The requirement for a sequence of significant results was used as a simple way to only take stable changes. These were selected manually on account of low individual-sensor noise (see results, between-sensor was much larger), although these parameters could be optimised in future work.

\subsection{Noise characterisation}

Sensors are subject to a range of noise sources. We performed an experiment to characterise: sensor noise per sample, which could vary per sensor, per-sensor bias, and per-placement bias, which accounts for variation such as due to variable packing, soil-moisture spatial heterogeneity, or insertion depth.

\hl{TODO: write up}

\begin{enumerate}
    \item Common pot with 40g of soil
    \item 4 sensors were placed in the soil, left to equilibriate, and measurements taken. This was repeated 8 times.
    \item Sensors were removed, 10 mL of water was added and the soil was re-mixed. Return to 2)
    \item By doing this, we were able to characterise short time-scale noise, potential sensor bias, and bias due to placement or packing.
\end{enumerate}

\subsection{Automated response curve calibration}

Consumer-grade soil moisture probes generally report coarse, normalized measurements. Here, we define soil water content as $Z$ as total water volume in a given container with fixed soil content (such as in a pot), instead of the usual $\theta$, which is the ratio of water to total soil volume. A well-described method in various contexts is to reconstruct a curve from added quantities when the quantity itself is not directly measurable, known in analytical chemistry as the method of standard additions \cite{ziegel2004statistics}. Similarly, here we assume that $Z$ itself is not directly measurable, although at a given watering, $\Delta Z$ is known, unlike $\Delta \theta$. We fix this for a given pot and soil volume, and refer to it as $W$. This is a simplifying assumption that disregards runoff and user measurement error.\\

For a given event at time $t_i$, and sensor $s$, we extract associated $X_{i,s}$. This leads to sampled pairs $(\Delta Z_{i,s}, X_{i,s})$. We asssume the relationship:

\[
X_{i,s} = h(Z_{i,s}) + \epsilon_i + \epsilon_s + \epsilon_p
\]

Where $\epsilon_i$ is measurement noise, $\epsilon_s$ noise offset per-sensor, and $\epsilon_p$ noise offset per-placement (accounting for factors such as sensor-specific offsets and variable depth and packing). In this equation, both the response function $h(\cdot)$ and the water content $Z$ are unknown. For simplicity, we assume $h$ is monotonically decreasing and refer to $g=h^{-1}$. In terms of the data, we have samples:

\[
\Big(\frac{\Delta Z}{\Delta X}\Big)_{i,s} \approx \Big(\frac{dz}{dx}\Big)_{i,s} + \epsilon_i + \epsilon_s + \epsilon_p
\]

where:

\[
\Big(\frac{dZ}{dX}\Big)_{i,s} = g'(X_{i,s})
\]

This is a basic regression problem except that it is required to estimate a derivative $g'$ instead of $g$ itself, and invert to get $h$. One complication is that for a given sample, we may like to use the midpoint time for a given estimate for the derivative: 

\[
\Big(\frac{\Delta Z}{\Delta X}\Big)_{i+1/2} = \frac{Z_{i+1} - Z{i}}{X_{i+1} - X{i}} = \frac{W}{X(Z + W) - X(Z)}
\]

The problem with this is that $X(Z + W)$ depends on noise attributable to spatial heterogeneity. In order to account for this, we introduce $M$, such that $X(Z + W\cdot M)$, where $M$ is non-negative and e.g. Gamma distributed (sensor values should decrease when water is added, but non-uniformly). We used:

\[
M \sim \Gamma(1/\sigma(z), \sigma(z))\]

with $\sigma$ linearly descreasing from $Z=0$ to $Z=Z_{max}$, in order to achieve $E[M] = 1$ and decreasing $var(M)$ consistent with observed variability descreasing with $Z$.  \\

We used Gaussian process regression \cite{williams2006gaussian} using Scikit-learn \cite{pedregosa2011scikit} to estimate $g'$. Then, we sample from $\hat{g'}$ and integrate to produce estimates and CIs for $\hat{g}$.

\[
g'(x) \sim \mathcal{GP}(m(X), k(X, X'))
\]

Estimation was performed with RBF kernel \cite{pedregosa2011scikit}. For integration, \texttt{scipy.integrate.cumulative\_trapezoid} \cite{virtanen2020scipy} was used. In practice, $X(Z_{max})$ is easiest endpoint to obtain because it is the value of $X$ at which all added water is lost as runoff. A prior linear calibration curve $g'_{prior}(x)$ was used by specifying constant derivative priors, as previously described \hl{cite}.

\begin{figure}[htb!]
  \centering
  \includegraphics[width=1.0\linewidth]{images/simulation_example.png}
  \caption{One sample of curve reconstruction process from simulated data. Top left: distribution of multiplicative noise as a function of $Z$ (95\% density). Top middle and top right: vertical (horizontal) lines indicate response $\Delta X$ on addition of $WM$ water. Middle-left and middle-middle: arrows indicate derivatives with error. Middle right and bottom left: plot of sampled derivatives. Bottom middle: GP regression of derivative curve $g'$. Bottom right: integrated estimated curve compared to real curve.}
  \label{fig:myfigure}
\end{figure}



\subsection{Validation}

\section{Results}

\subsection{Sensor noise}

TODO: sensor noise is mostly driven by packing and placement. Variation is not driven by sensor bias or by measurement error. Measurement error is low. 

\begin{figure}[htb!]
  \centering
  \includegraphics[width=1.0\linewidth]{images/equilibrium_noise_distribution.png}
  \caption{TODO: complete the data collection. Sensor measurements from real data. Top left: repeated sensor measurements in 3 time intervals corresponding to 0, 10, and 20 ml, $Z$ (or 0.0, 0.2, and 0.4 relative SWC, $\theta$). Top right: intra-placement variability, which was low. Bottom left: zero-centered total variability across intervals for each sensors, which was much greater than per-sample variability and \hl{TODO: and subject to test, may not be different between sensors}. Bottom right: X as a function of Z across moisture levels. Conclusions: variation due to placement contributes greatest of the sources considered \hl{TODO: statistical test}}.
  \label{fig:myfigure}
\end{figure}

\begin{enumerate}
    \item Sensor stationary noise and drying
    \item Per-sensor noise and per-insertion noise
\end{enumerate}

\subsection{Event detection}

\subsection{Calibration simulation}

\subsection{Validation}

\section{Discussion}

\section{Future work}

\bibliographystyle{plain}
\bibliography{references}

\end{document}
